COVID-19 pandemien er et bevis på, hvor komplekst det er at forberede sig og reagere på udbrud af smitsomme sygdomme. Ud fra et styringsmæssigt perspektiv krævede det koordinerede nationale og internationale indsatser. Ud fra et logistik perspektiv krævede det modstandsdygtige forsyningskæder for at distribuere medicinsk udstyr og vacciner. Ud fra et videnskabeligt perspektiv krævede det indsatser uden fortilfælde i størrelse og omfang til genomisk testning, sekvensering og modellering for at undersøge hvor virussen spredte sig og udviklede sig. Dog vil en omfattende reaktion mod sygdomsudbrud ikke afhænge af storstilet genomisk overvågning alene. Et godt eksempel er registerdata i Danmark, som var væsentlig for at afsløre regionale og demografiske mønstre af infektion. Under disse forudsætninger kan andre typer af data hjælpe at fremsætte værdifulde indsigter til udbrudsdynamik. Disse kan inkludere molekylære data fra biofysiske interaktioner mellem patogen og receptorer i værtsorganismer. En anden mulighed er data om økologi af naturlige reservoirer, som man kan opnå med ikke-invasive overvågningssystemer. Den slags datatyper er meget komplekse og det at udvikle effektive dataanalytiske metoder er altafgørende for at oversætte data til indsigter, som man kan basere afgørelser på inden for rimelige tidsrammer. Denne afhandling undersøger, hvordan vores forståelse af pandemier kan forbedres ved hjælp af nye datarepræsentationer af fylogeni, proteinstruktur og dyreadfærd.

Først udviklede jeg phylo2vec, en matematisk vektorrepræsentation af binære fylogenetiske træer. Designet med et særligt fokus på effektiv generering og datakompression muliggør denne kompakte repræsentation skalerbar analyse af fylogenetiske data med moderne \textit{machine learning} metoder. I det andet studie ønskede jeg efter at udforske en anden datastrøm: proteinstrukturer. Ved hjælp af en ny diskret repræsentation af proteinstruktur undersøgte jeg de evolutionære processer, der ligger bag affinitet mellem \textit{Betacoronavirus} Spike glykoproteiner og værtsreceptorer. I en første analyse af proteinsekvenser fandt jeg, at brugen af den humane ACE2-receptor er forbundet med genetisk rekombination i S1 underenheden af proteinet og med mutation i det N-terminale domæne af S1. Samtidig tyder strukturelle data på, at receptorskift medieres gennem gradvise strukturelle forbedringer inden for receptorbinding snarere end konformationel reorganisering. I den tredje studie undersøgte jeg potentialet for at detektere tidlige infektionssignaturer hos slagtekyllinger ved hjælp af \textit{deep learning}-baseret dyresporing. Ved at bruge \textit{Escherichia coli}, en betydelig trusel i den danske fjerkræbranche, som modelpatogen, fandt jeg beviser for at denne teknik kan bruges til at måle adfærdsmønstre, som er forbundet med forringet velfærd. Dette resultat viser kunstig intelligens' potentiale for at automatisere omdannelse af videoovervågningsdata til kvantitative dyrevelfærdsmålinger og, alt i alt, drive effektive løsninger til epidemiovervågning i fjerkræpopulationer.

Samlet set illustrerer denne afhandling den iboende tværfaglige karakter af et beredskab for pandemier. Den viser, at forskellige datastrømme kan integreres i sundhedsinfrastrukturer via avanceret datateknik. Ud over konventionel genetisk og videobaseret overvågning vil nye analytiske metoder accelerere detektion og analyse af patogener og uddybe vores forståelse af nye trusler fra infektionssygdomme.

% In the third data stream, I investigated the feasibility of detecting early infection signatures in broiler chickens using deep learning-based animal tracking. Using \textit{Escherichia coli}, a major threat in poultry farming in Denmark, as a model pathogen, I found evidence that simple behavioural patterns tracked from videography are associated with impaired welfare. This result showcases the potential of artificial intelligence to power low-cost solutions for epidemic monitoring in poultry populations, enabling the automated transformation of video surveillance into quantitative animal welfare metrics.

% First, I developed phylo2vec, a mathematical integer vector representation of binary phylogenetic trees. Designed with efficient sampling and data storage in mind, this compact representation paves the way for scalable analysis of phylogenetic data using modern machine learning techniques. Exploring a second data stream, I used protein structures to investigate the evolutionary mechanisms underlying the affinity of \textit{Betacoronavirus} spike proteins to host receptors. By casting protein structures into a discrete representation, I found that the use of the human ACE2 receptor is linked with molecular recombination in the gene S1 subunit and with amino acid substitution rates in the N-terminal domain. Meanwhile, structural data suggest that receptor switching is mediated through gradual structural refinements within the receptor-binding domain rather than large-scale conformational reorganisation. In the third data stream, I investigated the feasibility of detecting early infection signatures in broiler chickens using deep learning-based animal tracking. Using \textit{Escherichia coli}, a major threat in poultry farming in Denmark, as a model pathogen, I found evidence that simple behavioural patterns tracked from videography are associated with impaired welfare. This result showcases the potential of artificial intelligence to power low-cost solutions for epidemic monitoring in poultry populations, enabling the automated transformation of video surveillance into quantitative animal welfare metrics.

% Ved at udforske en anden datastrøm har jeg omdannet proteinstrukturer til diskrete repræsentationer og brugt disse til at undersøge de evolutionære mekanismer, der ligger bag affinitet mellen antigener og værtsreceptorer. I en analyse af sekvensdata på aminosyrer af \textit{Betacoronavirus} spike glycoproteiner fandt jeg, at brugen af den humane ACE2-receptor er forbundet med molekylær rekombination i underenheden til genet S1 og med aminosyresubstitutionshastigheder i det N-terminale domæne. Samtidig tyder strukturelle data på, at receptorskift medieres gennem gradvise strukturelle forbedringer inden for receptorbinding snarere end storstilet konformationel reorganisering. I den tredje datastrøm undersøgte jeg muligheden for deep learning-baseret sporing af dyr for at detektere tidlige infektionssignaturer hos slagtekyllinger. Ved at bruge \textit{Escherichia coli}, en stor trussel i fjerkræavl i Danmark, som modelpatogen, fandt jeg beviser fra videooptagelser for, at simple adfærdsmønstre er forbundet med forringet velfærd. Dette placerer derfor computer vision som en kraftfuld og billig løsning til epidemiovervågning i fjerkræpopulationer via automatiseret omdannelse af videoovervågning til kvantitative adfærdsmålinger.

% Samlet set illustrerer denne afhandling den iboende tværfaglige karakter af et beredskab for pandemier. Den viser, at forskellige datastrømme kan integreres i sundhedsinfrastrukturer via avanceret datateknik. Ud over konventionel genetisk og videobaseret overvågning vil nye analytiske metoder accelerere detektion og analyse af patogener og uddybe vores forståelse af nye trusler fra infektionssygdomme.



% % Et beredskab for pandemier er en vigtig byggesten i den globale sundhedsinfrastruktur. Ud over logistiske og forvaltningsmæssige udfordringer kræver et robust beredskab en forskningsindsats på tværs af flere videnskabelige discipliner. Den seneste COVID-19-pandemi eksemplificerer denne kompleksitet og viser, at en omfattende respons på sygdomsudbrud muligvis ikke kun afhænger af omfattende testning for infektion. I stedet kan vi udvide vores forståelse af sygdomsudbrud betydeligt ved at inkludere komplementære datastrømme. Disse datastrømme kan omfatte data om udvikling af patogener på tværs af millioner af patienter, biofysiske interaktioner mellem patogener og værtsmolekyler samt ikke-invasive detektionssystemer i dyrereservoirer. Disse datatyper er yderst komplekse og udvikling af beregningsmæssigt effektive analysemetoder er afgørende for at omsætte data til brugbar indsigt inden for operationelt relevante tidsrammer. Denne afhandling undersøger, hvordan vores forståelse af pandemier kan forbedres ved hjælp af nye datarepræsentationer af fylogeni, proteinstruktur og dyreadfærd.


% \textbf{Things I am unsure of:
% }
% % - data streams: datastrømme? Not in the dictionary, but websites use it. https://ordnet.dk/ddo/ordbog?query=datastrømmen

% % - integer... the word heltalsvektorrepræsentation does not seem to exist even though google translate wants it to. I propose something without the interger part maybe? Or have it as a seperate word before. Now I have written "matematisk vektorrepræsentation af binære fylogenetiske træer" which is does not cover the interger part.

% % - discrete representation: diskrete repræsentationer? Should be fine but maybe ask Frederik? I cannot find anything on google on this in Danish Math things.

% % - governing antigen affinity to host receptors:  styrer affinitet mellen antigener og værtsreceptorer. Not 100 procent sure about this translation because of my lack of biology knowlegde.

% - amino acid substitution rates: aminosyresubstitutionshastigheder. This is the direct translation but does not seem to exist. Could be "substitionshastigheder for aminosyre".

% - data engeneering: datatektik. One of those words that might be better in english

% - also double check this sentence with someone with biology knowlegde: I en analyse af sekvensdata på aminosyrer af Betacoronavirus spike glycoproteiner fandt jeg, at brugen af den humane ACE2-receptor er forbundet med molekylær rekombination i underenheden til genet S1 og med aminosyresubstitutionshastigheder i det N-terminale domæne

% - last sentence: is it new treats or new diseases? The "nye" might need to be moved